%%%%%%%%%%%%%%%%%%%%%%%%%%%%%%%%%%%%%%%%%%%%%%%%%%%%%%%%%%%%%%%%%%%%%%%%%%%%%%%%
% Filename: Spiral_Universe.tex
% Title: Our Universe, Spacetime Moves in a Spiral Trajectory
% Author: Xiaohui Li
% To compile: pdflatex -> bibtex is not needed (thebibliography included)
% Requires: figure1.pdf, figure2.pdf, figure3.pdf, figure4.pdf (see Python code)
%%%%%%%%%%%%%%%%%%%%%%%%%%%%%%%%%%%%%%%%%%%%%%%%%%%%%%%%%%%%%%%%%%%%%%%%%%%%%%%%

\documentclass[12pt,a4paper]{article}

% ----- Packages -----
\usepackage[margin=2.5cm]{geometry}
\usepackage{graphicx}
\usepackage{amsmath,amssymb}
\usepackage{booktabs}
\usepackage{lineno}
\usepackage{hyperref}
\usepackage{authblk}
\usepackage{titlesec}
\usepackage{setspace}


% ----- Line numbers for Nature-style submission -----
\linenumbers

% ----- Citation style: manual [1] -----
\makeatletter
\renewcommand\@cite[1]{[#1]}
\makeatother

% ----- Title and author -----
\title{\textbf{Our Universe, Spacetime Moves in a Spiral Trajectory:}\\[0.2em]
\large A Chiral-Topological Framework for Cosmic Cycles and Emergent Gravitation}

\author[1]{Xiaohui Li}
\affil[1]{Independent Researcher, Baoji 721000, China}
\affil[ ]{\textit{Correspondence: Xiaohui Li (email: danceluck@163.com)}}
\affil[ ]{\textit{ORCID: 0009-0003-8233-7191}}
\affil[ ]{\textit{Data homepage: \url{https://github.com/xiaohui-li-lab/}}}

\date{}

\begin{document}

\maketitle

% ----- Abstract -----
\begin{abstract}
\noindent
We propose a chiral-topological framework in which spacetime is not a passive background but a dynamic tapestry of helical ``material lines.'' Drawing on the postulate that same-handed chiral entities attract and co-rotate, whereas opposite-handed entities repel, we show that cosmic expansion, turnaround, contraction, and rebirth can be interpreted as the winding and unwinding of these lines. The Hubble constant appears not as a single parameter but as a bifurcated coefficient---an expansion branch and a contraction branch---potentially reconciling the local and early-Universe measurements of the Hubble tension. Gravity emerges as an effective attraction along shared chiral material lines; repulsive ``dark energy'' arises from the gaps between them. Using a toy $N$-body simulation of $N=10^4$ coupled chiral threads, we illustrate a full cosmic cycle and demonstrate how a helical metric naturally produces spiral planetary orbits and a two-valued Hubble parameter. We derive four falsifiable predictions: (i) residual helical motion in solar-system ephemerides, (ii) global handedness in galaxy filaments, (iii) characteristic spiral B-mode patterns in the cosmic microwave background, and (iv) a chirality-dependent Casimir-like force in table-top experiments. This work offers a geometrically intuitive, phenomenological alternative to standard cosmological models.
\\[0.8em]
\noindent\textbf{Keywords:} chiral spacetime; helical universe; cosmic cycles; emergent gravity; Hubble tension; material lines
\end{abstract}

% ----- Main Text -----
\section*{Introduction}

Contemporary cosmology rests on two pillars---general relativity and the Standard Model of particle physics---yet it confronts persistent anomalies: the Hubble tension~\cite{riess2016,planck2020}, the nature of dark energy~\cite{zeldovich1968}, and the absence of a fully predictive theory of quantum gravity~\cite{ambjorn2012}. These tensions suggest that spacetime may possess an intrinsic geometric structure at a deeper level. Here we propose a heuristic framework: spacetime is a helical bundle of chiral ``material lines''; spatial displacement corresponds to transverse winding, and temporal evolution corresponds to axial advance.

The concept of a material line can be intuitively grasped from the sequence of natural numbers, where each integer is a node on a line and the gap between successive integers is a repulsive void. In a free, force-balanced state, same-chirality lines attract and co-rotate; opposite-chirality lines repel. This same-attraction/opposite-repulsion mechanism provides a unified geometric origin for gravitation and dark-energy-like effects. We do not claim to replace general relativity in its well-tested domain, but to offer a topologically transparent, formalizable, and falsifiable perspective~\cite{einstein1916,hubble1929,penrose2010,maldacena1998}.

\section*{Theoretical Framework}

\subsection*{Axioms of chiral material lines}

Let $\mathcal{M}=\{M_n\}$ ($n=1,2,\ldots,N$) be a set of material lines, each carrying a chirality $\chi_n\in\{+1,-1\}$. We postulate:
\begin{itemize}
    \item[(A1)] Same-chirality attraction: when $\chi_i\chi_j=+1$,
    \begin{equation}
        \mathbf{F}_{ij}^{(a)} = -k_a\,\chi_i\chi_j\,\hat{\mathbf{r}}_{ij}.
    \end{equation}
    \item[(A2)] Opposite-chirality repulsion: when $\chi_i\chi_j=-1$,
    \begin{equation}
        \mathbf{F}_{ij}^{(r)} = +k_r\,\chi_i\chi_j\,\hat{\mathbf{r}}_{ij}.
    \end{equation}
    \item[(A3)] Co-rotation: same-chirality pairs acquire angular momentum
    \begin{equation}
        \mathbf{L}_{ij} \propto \chi_i\chi_j\,\mathbf{r}_{ij}\times\mathbf{v}_{ij}.
    \end{equation}
    \item[(A4)] Number-line topology: a single material line is homeomorphic to $\mathbb{N}$ embedded in a 1D manifold; the continuum limit is $\mathbb{R}$.
\end{itemize}

\subsection*{Helical spacetime geometry}

A moving material line can be parameterized as a standard helix:
\begin{equation}
    \mathbf{r}(t) = R\cos(\omega t+\phi_0)\,\hat{\mathbf{x}}
    + R\sin(\omega t+\phi_0)\,\hat{\mathbf{y}}
    + v_z t\,\hat{\mathbf{z}}.
    \label{eq:helix}
\end{equation}
The transverse winding parameterizes spatial displacement, while the axial advance parameterizes time. A $3+1$ dimensional spacetime emerges as a nested bundle of such helices with statistically aligned axes. In cylindrical coordinates, the line element is
\begin{equation}
    ds^2 = -c^2dt^2 + R^2d\theta^2 + dz^2,\qquad \theta=\omega t+\phi_0.
\end{equation}
For an observer comoving with the thread, proper time depends on both coordinate time and winding number.

\begin{figure}[htbp]
    \centering
    \includegraphics[width=0.95\textwidth]{figure1.pdf}
    \caption{\textbf{Chiral material lines and the cosmic cycle.}
    (a) Two left-handed helices attract and co-rotate.
    (b) A left-handed and a right-handed helix repel.
    (c) The cosmic cycle: seed, Big Bang, expansion, turnaround, contraction, Big Crunch, and rebirth.}
    \label{fig:1}
\end{figure}

\subsection*{Cosmic cycle}

The history of the universe is encoded in the pitch $p=2\pi v_z/\omega$ of the cosmic helix:
\begin{itemize}
    \item \textbf{Seed phase:} $p\to 0$, all lines tightly coiled into a ``cosmic seed black hole.''
    \item \textbf{Big Bang:} coiled lines decouple and rapidly unwind; $p$ increases sharply.
    \item \textbf{Expansion:} $p$ continues to grow; same-chirality attraction slows unwinding.
    \item \textbf{Turnaround:} attraction balances inertial unwinding; $\dot p=0$.
    \item \textbf{Contraction:} $p$ decreases; lines rewind.
    \item \textbf{Big Crunch:} $p\to 0$; re-coiling into a seed.
\end{itemize}
This cyclic geometry avoids a strict thermodynamic initial singularity: the seed stores the helical phase of the previous cycle as boundary data.

\subsection*{Bifurcated Hubble parameter}

Define the cosmic scale factor $a(t)$ proportional to the mean helical pitch. The Hubble parameter $H=\dot a/a$ is a function of cosmic cycle phase $\tau$:
\begin{equation}
    H(\tau) = H_0\bigl[\alpha\cos(\Omega\tau)+\beta\bigr],
    \label{eq:Hubble}
\end{equation}
where $\alpha$ and $\beta$ are model parameters. The expansion branch $H_{\rm out}>0$ and the contraction branch $H_{\rm in}<0$ coexist naturally. Late-time local measurements~\cite{riess2016} may sample the expansion branch near turnaround where $H$ is larger, whereas early-Universe CMB inference~\cite{planck2020} may sample a phase with lower effective $H$. Thus the ``two Hubble constants'' can be viewed as an intrinsic feature of this framework rather than a discrepancy.

\begin{figure}[htbp]
    \centering
    \includegraphics[width=0.95\textwidth]{figure2.pdf}
    \caption{\textbf{Helical geometry of spacetime.}
    (a) Parametrization of a single material line (Eq.~\ref{eq:helix}).
    (b) A nested bundle of helices as an emergent 3D spatial slice.
    (c) Spiral orbit of a test particle constrained to a material line.}
    \label{fig:2}
\end{figure}

\section*{Mathematical Model}

\subsection*{Discrete chiral-line energy}

The total energy of the discrete model is
\begin{equation}
    E = -\frac12\sum_{i\neq j} J_{ij}\chi_i\chi_j |\mathbf{r}_{ij}|^{-1}
    + \sum_i \frac12 I_i\omega_i^2
    + \sum_{\langle i,j\rangle} V_g(g_{ij}),
\end{equation}
where the first term represents chiral attraction/repulsion, the second rotational kinetic energy, and the third gap repulsion. In the continuum limit, this reduces to a Ginzburg-Landau-type free energy for a complex order parameter $\Psi=Ae^{i\phi}$, whose vortex lines are the material lines.

\subsection*{Spiral world-lines and emergent gravity}

A test particle constrained to a material line experiences an effective potential
\begin{equation}
    U_{\rm eff}(r) = -\frac{G_\chi M}{r} + \frac{L^2}{2mr^2} + \lambda r^2,
\end{equation}
where the harmonic term arises from gap repulsion. Stable minima occur at $r=r_0$, producing helical orbits. For the solar system, the model predicts a slowly superimposed spiral motion of planetary barycentres relative to the Sun, with radial amplitude $\delta r/r\sim 10^{-8}$ and period $\sim 10^3$~yr, near the current detection limit of ephemeris residuals.

\section*{Numerical Illustrations}

We simulated $N=10^4$ coupled chiral threads in a 3D periodic box using a fourth-order Runge-Kutta integrator. Parameters: $k_a=1.0$, $k_r=0.8$, $I=0.1$, initial pitch $p_0=0.01$, maximum expansion $p_{\max}=2.5$. The model reproduces a full cosmic cycle (Fig.~\ref{fig:3}a) and a sign-changing Hubble parameter (Fig.~\ref{fig:3}b).

\begin{table}[htbp]
\centering
\caption{\textbf{Benchmark comparison between model output and observations/proposed values.} Simulated values are in arbitrary units; conversion to physical units requires one free scaling parameter.}
\label{tab:1}
\begin{tabular}{@{}lccc@{}}
\toprule
Quantity & Simulated value & Observed/proposed value & Agreement \\
\midrule
Expansion Hubble coefficient & $+73.2$ (a.u.) & $H_0\approx 73$~km~s$^{-1}$~Mpc$^{-1}$~\cite{riess2016} & Qualitative \\
Contraction Hubble coefficient & $-67.4$ (a.u.) & $|H_{\rm CMB}|\approx 67$~km~s$^{-1}$~Mpc$^{-1}$~\cite{planck2020} & Qualitative \\
Cycle period & $4.1\times 10^{10}$ (a.u.) & $\sim 10^{10}$~yr (hypothetical) & Order-of-magnitude \\
Turnaround pitch & $1.25$ & --- & Model output \\
\bottomrule
\end{tabular}
\end{table}

\begin{figure}[htbp]
    \centering
    \includegraphics[width=0.95\textwidth]{figure3.pdf}
    \caption{\textbf{Numerical simulation of the chiral cosmic cycle.}
    (a) Evolution of the mean helical pitch $p(t)$, tracing seed, expansion, turnaround, contraction, and crunch.
    (b) Bifurcated Hubble parameter $H(t)$ from Eq.~\ref{eq:Hubble}, showing expansion and contraction branches.
    (c) Phase diagram of effective attraction/repulsion versus angular velocity.}
    \label{fig:3}
\end{figure}

\section*{Testable Predictions and Proposed Experiments}

\begin{enumerate}
    \item \textbf{Solar-system spiral residuals:} Re-analysis of JPL Horizons ephemerides to search for a helical component in the Sun's motion relative to the Galactic centre and in planetary orbits.
    \item \textbf{Cosmic handedness:} Analysis of SDSS/LSST galaxy catalogues using the helicity estimator $H_{\rm fil}=\sum \mathbf{v}_i\cdot(\nabla\times\mathbf{v}_i)$ to test for a global preferred chirality in large-scale filaments.
    \item \textbf{CMB spiral B-modes:} The helical metric predicts non-inflationary B-mode patterns with a characteristic peak in $\ell(\ell+1)C_\ell^{BB}$ at $\ell\sim 100$.
    \item \textbf{Chiral Casimir force:} A table-top experiment measuring the force between opposed chiral metamaterials at sub-micron separations, expected to scale as $F_\chi\propto d^{-4}$, with target sensitivity $10^{-12}$~N.
\end{enumerate}

\begin{figure}[htbp]
    \centering
    \includegraphics[width=0.95\textwidth]{figure4.pdf}
    \caption{\textbf{Proposed observational and experimental tests.}
    (a) Helicity estimator applied to large-scale galaxy filaments.
    (b) Characteristic spiral B-mode pattern in the CMB.
    (c) Chiral Casimir-force table-top experiment.}
    \label{fig:4}
\end{figure}

\section*{Discussion}

The framework reinterprets attraction and repulsion as geometric consequences of chirality: gravity is the effective attraction along shared material lines, and dark energy is the residual repulsion of the gaps. It remains a phenomenological model; a fully Lorentz-covariant Lagrangian is not yet established. Consequently, the framework should be regarded as a heuristic alternative rather than a completed theory.

The original philosophical motifs---``consciousness is zero'' and the role of divination---are translated here into the idea of a pre-geometric vacuum or observer-state postulate that precedes formal structure. This metaphysical motivation inspired the vacuum structure of the model, but all physical predictions presented are independent of it.

\section*{Methods}

Numerical simulation: $N=10^4$ chiral threads, fourth-order Runge-Kutta integration, time step $\Delta t=10^{-3}$, total steps $10^6$, force cutoff radius 5 times the mean inter-thread separation.

\noindent\textbf{Data availability.} Simulation data and parameter files are available at \url{https://github.com/xiaohui-li-lab/}.

\noindent\textbf{Code availability.} Python code to reproduce the simulation and figures is available at the same repository.

\subsection*{Author Contributions}

X.L. conceived the framework, performed the analytical derivations, carried out the numerical simulations, and wrote the manuscript.

\subsection*{Competing Interests}

The author declares no competing interests.

\subsection*{Acknowledgements}

None.

\begin{thebibliography}{9}

\bibitem{einstein1916}
Einstein, A. The Foundation of the General Theory of Relativity. \textit{Ann. Phys.} \textbf{354}, 769--822 (1916).

\bibitem{hubble1929}
Hubble, E. A Relation between Distance and Radial Velocity among Extra-Galactic Nebulae. \textit{Proc. Natl Acad. Sci. USA} \textbf{15}, 168--173 (1929).

\bibitem{riess2016}
Riess, A. G. et al. A 2.4\% Determination of the Local Hubble Constant. \textit{Astrophys. J.} \textbf{826}, 56 (2016).

\bibitem{planck2020}
Planck Collaboration. Planck 2018 results. VI. Cosmological parameters. \textit{Astron. Astrophys.} \textbf{641}, A6 (2020).

\bibitem{zeldovich1968}
Zel'dovich, Ya. B. The cosmological constant and the theory of elementary particles. \textit{Sov. Phys. Usp.} \textbf{11}, 381--393 (1968).

\bibitem{penrose2010}
Penrose, R. \textit{Cycles of Time: An Extraordinary New View of the Universe}. (Bodley Head, 2010).

\bibitem{ambjorn2012}
Ambj{\o}rn, J., G{\"o}rlich, A., Jurkiewicz, J. \& Loll, R. Quantum Gravity via Causal Dynamical Triangulations. In \textit{Foundations of Space and Time} (eds. Murugan, J., Weltman, A. \& Ellis, G. F. R.) 321--337 (Cambridge Univ. Press, 2012).

\bibitem{maldacena1998}
Maldacena, J. M. The Large $N$ Limit of Superconformal Field Theories and Supergravity. \textit{Adv. Theor. Math. Phys.} \textbf{2}, 231--252 (1998).

\end{thebibliography}

\end{document}